\begin{abstract}
Most diffusion and flow-matching generators define the prior, probability path, and prediction target in the same representation space. Latent diffusion improves efficiency by moving this path into an autoencoder latent space, but the final sample is still produced by a separately trained decoder. This separation creates a mismatch: the generator is optimized for latent-space prediction, while final quality depends on how the decoder handles generated latents that may differ from clean encoder outputs. We introduce CrossFlow, a cross-space flow formulation that maps noisy latent inputs directly to pixel-space images. The key technical step is a velocity-free one-step objective: the latent trajectory defines the training path, but the supervised prediction is an image rather than a latent displacement. This lets one model act both as a one-step latent-to-pixel generator and as a decoder replacement for latent diffusion pipelines. On class-conditional ImageNet-1k at $256\times256$, CrossFlow-XL achieves 1.62 FID with one function evaluation. Ablations show that the latent encoder and pixel-space perceptual and adversarial losses are important for fidelity. These results indicate that cross-space flow objectives can combine the efficiency of latent representations with direct pixel-space supervision, without requiring a separate decoder at inference.
\end{abstract}
